% ===== PRÉAMBULE CANONIQUE — à coller VERBATIM en tête de CHAQUE .tex =====
\documentclass[11pt,a4paper]{article}
\usepackage[utf8]{inputenc}\usepackage[T1]{fontenc}\usepackage{lmodern}
\usepackage[french]{babel}
\usepackage{amsmath,amssymb,mathtools}\usepackage{icomma}
\usepackage{siunitx}\sisetup{locale=FR}  % \qty{}{}, \si{}, \ang{}
\usepackage{xcolor}\usepackage{colortbl}
\usepackage{tikz}\usepackage{pgfplots}\pgfplotsset{compat=1.18}
\usepackage[siunitx,europeanresistors]{circuitikz} % circuits (élec.)
\usepackage[most]{tcolorbox}\usepackage{enumitem}\usepackage{array,booktabs}
\usepackage[a4paper,margin=2.2cm]{geometry}
\usetikzlibrary{arrows.meta,calc,angles,quotes,positioning,patterns,%
  backgrounds,shapes.geometric,decorations.pathmorphing,decorations.markings,intersections}
\definecolor{primary}{RGB}{222,49,99}\definecolor{primarydark}{RGB}{150,22,55}
\definecolor{defcol}{RGB}{0,114,178}\definecolor{methcol}{RGB}{52,92,148}
\definecolor{excol}{RGB}{0,135,90}\definecolor{attncol}{RGB}{200,40,55}
\tcbset{boxbase/.style={enhanced,breakable,boxrule=0.9pt,arc=2.5pt,
  left=3mm,right=3mm,top=2.3mm,bottom=2.3mm,fonttitle=\bfseries,coltitle=white}}
% --- boîtes TUTORIEL (noms EXACTS, ne pas renommer) ---
\newtcolorbox{cle}[1][]{boxbase,colback=defcol!6,colframe=defcol,
  title={\textsc{À retenir}\ifx\relax#1\relax\relax\else\textmd{ : \itshape #1}\fi}}
\newtcolorbox{methode}[1][]{boxbase,colback=methcol!7,colframe=methcol,sharp corners,
  title={\textsc{Méthode}\ifx\relax#1\relax\relax\else\textmd{ --- \itshape #1}\fi}}
\newtcolorbox{exemplebox}[1][]{boxbase,colback=excol!5,colframe=excol,
  title={\textsc{Exemple}\ifx\relax#1\relax\relax\else\textmd{ : \itshape #1}\fi}}
\newtcolorbox{piege}[1][]{boxbase,colback=attncol!6,colframe=attncol,
  title={\textsc{Piège}\ifx\relax#1\relax\relax\else\textmd{ : \itshape #1}\fi}}
\newtcolorbox{remarque}[1][]{boxbase,colback=black!4,colframe=black!45,
  title={\textsc{Remarque}\ifx\relax#1\relax\relax\else\textmd{ : \itshape #1}\fi}}
% --- boîtes EXERCICE / CORRIGÉ (noms EXACTS, auto-comptées) ---
\newtcolorbox[auto counter]{exo}[1][]{boxbase,colback=primary!4,colframe=primary,
  title={\textsc{Exercice}~\thetcbcounter\ifx\relax#1\relax\relax\else\textmd{ --- \itshape #1}\fi}}
\newtcolorbox[auto counter]{corrige}[1][]{boxbase,colback=excol!4,colframe=excol,
  title={\textsc{Corrigé}~\thetcbcounter\ifx\relax#1\relax\relax\else\textmd{ --- \itshape #1}\fi}}
\newcommand{\vv}[1]{\vec{#1}}\newcommand{\ds}{\displaystyle}
\newcommand{\R}{\mathbb{R}}\newcommand{\N}{\mathbb{N}}\newcommand{\Z}{\mathbb{Z}}
% =========================================================================
\begin{document}
\sisetup{per-mode=symbol}
% ================= FACILE =================

\begin{exo}[calcul direct]
Deux sphères de plomb, de masses $m_A=\qty{10}{\kilogram}$ et
$m_B=\qty{10}{\kilogram}$, ont leurs centres distants de $d=\qty{0.50}{\metre}$.
Calcule l'intensité de la force de gravitation qu'elles exercent l'une sur l'autre.
On donne $\mathcal{G}=\num{6.67e-11}\ \si{\newton\metre\squared\per\kilogram\squared}$.
\end{exo}

\begin{exo}[les deux forces]
L'intensité de la force exercée par $A$ sur $B$ vaut $F_{A/B}=\qty{5.0e-8}{\newton}$.
\begin{enumerate}[label=\alph*)]
  \item Quelle est l'intensité de la force exercée par $B$ sur $A$ ?
  \item Ces deux forces ont-elles le même sens ou des sens opposés ?
  \item Sur quelle droite sont-elles portées ?
\end{enumerate}
\end{exo}

\begin{exo}[l'effet des masses]
Deux corps s'attirent avec une force d'intensité $F_0$. Sans changer la distance,
exprime la nouvelle force en fonction de $F_0$ si :
\begin{enumerate}[label=\alph*)]
  \item on double la masse $m_A$ ;
  \item on triple la masse $m_B$ ;
  \item on double les deux masses en même temps.
\end{enumerate}
\end{exo}

\begin{exo}[l'effet de la distance]
Deux corps s'attirent avec une force d'intensité $F_0$. Sans changer les masses,
exprime la nouvelle force en fonction de $F_0$ si :
\begin{enumerate}[label=\alph*)]
  \item on double la distance $d$ ;
  \item on triple la distance $d$ ;
  \item on divise la distance $d$ par $2$.
\end{enumerate}
\end{exo}

\begin{exo}[force Terre--satellite]
Un satellite de masse $m=\qty{500}{\kilogram}$ a son centre à la distance
$d=\qty{8.0e6}{\metre}$ du centre de la Terre ($M_T=\qty{6.0e24}{\kilogram}$).
Calcule l'intensité de la force de gravitation exercée par la Terre sur le satellite.
\end{exo}

\end{document}
